% Figura -- Fluxo de implementação do localizador como firmware.
%
% Compilar com:  pdflatex fig-pipeline.tex
% Gera fig-pipeline.pdf, incluído no corpo do relatório.
%
% Mesmas convenções e paleta do fluxograma da monografia (fluxograma-metodologia.tex):
%   - seta simples entre caixas ........ etapa seguinte opera sobre a saída da anterior;
%   - barra de bifurcação / junção ..... ramos operam sobre a MESMA entrada e são depois
%                                        combinados (no microcontrolador, executados em
%                                        sequência, por haver um único núcleo).

\documentclass[border=4pt]{standalone}

\usepackage[T1]{fontenc}
\usepackage[utf8]{inputenc}
\usepackage[brazil]{babel}
\usepackage{lmodern}
\usepackage{tikz}
\usetikzlibrary{arrows.meta}

\definecolor{cxfill}{HTML}{BEE3F0}
\definecolor{cxline}{HTML}{1F5C7A}
\definecolor{bdfill}{HTML}{F4F9FC}
\definecolor{bdline}{HTML}{8FB3C6}
\definecolor{rsfill}{HTML}{7FC3DC}
\definecolor{ntcolor}{HTML}{3C6A82}

\begin{document}
\sffamily
\hyphenpenalty=10000 \exhyphenpenalty=10000
\begin{tikzpicture}[
  x=1cm, y=1cm,
  etapa/.style={
    rounded corners=2pt, draw=cxline, fill=cxfill, line width=0.6pt,
    text width=#1, align=center, inner sep=3pt, minimum height=1.10cm,
    font=\scriptsize},
  etapa/.default=3.20cm,
  resultado/.style={
    rounded corners=3pt, draw=cxline, fill=rsfill, line width=1.1pt,
    text width=#1, align=center, inner sep=4pt, minimum height=1.10cm,
    font=\scriptsize\bfseries},
  resultado/.default=3.20cm,
  faixa/.style={rounded corners=5pt, draw=bdline, fill=bdfill, line width=0.7pt},
  titulo/.style={font=\small\bfseries, text=cxline},
  seta/.style={-{Stealth[length=4pt,width=4pt]}, draw=cxline, line width=0.7pt},
  liga/.style={draw=cxline, line width=0.7pt},
  setona/.style={-{Stealth[length=8pt,width=8pt]}, draw=cxline, line width=2pt},
  tracejada/.style={-{Stealth[length=4pt,width=4pt]}, draw=cxline, line width=0.7pt,
                    dash pattern=on 3pt off 2pt},
  barra/.style={draw=cxline, line width=1.6pt},
  nota/.style={font=\scriptsize\itshape, text=ntcolor, inner sep=1.5pt},
  par/.style={font=\scriptsize\itshape, text=cxline, inner sep=1.5pt},
]

% ===================== FAIXA 1 -- preparação em Python =====================
\fill[faixa] (0.2,-0.2) rectangle (16.2,-2.75);
\node[titulo, anchor=west] at (0.55,-0.62) {1.\ \ Preparação em \textit{Python} (computador)};

\node[etapa] (a1) at (2.2,-1.85) {Base de atributos da monografia, a 40\,dB};
\node[etapa] (a2) at (6.1,-1.85) {Partição fixa por solo e seleção sobre o treinamento};
\node[etapa] (a3) at (10.0,-1.85) {Treinamento dos modelos: RL, AD, FA e RNA};
\node[etapa] (a4) at (13.9,-1.85) {Geração de código C: modelo, configuração e janelas de teste};

\draw[seta] (a1) -- (a2);
\draw[seta] (a2) -- (a3);
\draw[seta] (a3) -- (a4);
\node[nota, above] at (2.2,-1.22) {7990 simulações, 7056 atributos};
\node[nota, above] at (6.1,-1.22) {6345 treino / 1645 teste};

\draw[setona] (13.9,-2.75) -- (13.9,-3.40);
\node[nota, anchor=west] at (14.15,-3.05) {código gerado};

% ===================== FAIXA 2 -- firmware ===================
\fill[faixa] (0.2,-3.40) rectangle (16.2,-9.25);
\node[titulo, anchor=west] at (0.55,-3.82) {2.\ \ \textit{Firmware} em C (Cortex-M4F simulado no Renode)};

\node[etapa=3.6cm] (b1) at (2.6,-5.00)
  {Janela bruta: 7 sinais da subestação, 10 ciclos pré e pós-falta};
\node[nota, below] at (2.6,-5.62) {2 $\times$ 7 $\times$ 2560 amostras};

% bifurcação: as três transformadas operam sobre a MESMA janela
\draw[liga] (b1.east) -- (5.6,-5.00) -- (5.6,-5.45);
\draw[barra] (5.6,-5.45) -- (12.0,-5.45);
\node[par, anchor=west] at (6.7,-5.26) {sobre a mesma janela};

\node[etapa=2.45cm] (b2) at (6.6,-6.35) {TDF ciclo a ciclo (CMSIS-DSP, 256 pontos)};
\node[etapa=2.45cm] (b3) at (9.3,-6.35) {TS nas ordens harmônicas (10 TRF de 256 pontos)};
\node[etapa=2.45cm] (b4) at (12.0,-6.35) {TW não decimada (db4, 8 níveis)};
\draw[seta] (6.6,-5.45) -- (b2.north);
\draw[seta] (9.3,-5.45) -- (b3.north);
\draw[seta] (12.0,-5.45) -- (b4.north);

% junção
\draw[liga] (b2.south) -- (6.6,-7.20);
\draw[liga] (b3.south) -- (9.3,-7.20);
\draw[liga] (b4.south) -- (12.0,-7.20);
\draw[barra] (6.6,-7.20) -- (12.0,-7.20);

\node[etapa=3.6cm] (b5) at (2.6,-8.35) {Parâmetros elétricos e métricas estatísticas};
\draw[seta] (9.3,-7.20) -- (9.3,-7.45) -- (2.6,-7.45) -- (b5.north);
\node[etapa=2.8cm] (b6) at (7.3,-8.35) {Modelo gerado};
\node[resultado=3.4cm] (b7) at (12.4,-8.35) {Região estimada (1 a 4)};
\draw[seta] (b5) -- (b6);
\draw[seta] (b6) -- (b7);

\draw[setona] (8.2,-9.25) -- (8.2,-9.90);

% ===================== FAIXA 3 -- verificação e medição ===================
\fill[faixa] (0.2,-9.90) rectangle (16.2,-12.30);
\node[titulo, anchor=west] at (0.55,-10.32) {3.\ \ Verificação e medição};

\node[etapa=3.9cm] (c1) at (3.0,-11.45)
  {Comparação de atributos e predições com a referência em \textit{Python}};
\node[etapa=3.9cm] (c2) at (8.2,-11.45)
  {Concordância mínima de 99\,\% nas 120 janelas de teste};
\node[resultado=3.9cm] (c3) at (13.4,-11.45)
  {Medição de flash, RAM, pilha e instruções executadas};
\draw[seta] (c1) -- (c2);
\draw[seta] (c2) -- (c3);

\end{tikzpicture}
\end{document}
